\section{Factorización, relojes aritméticos y jerarquía de trazas}
\label{lectura:manuscrito-02}
El producto nativo permite distinguir dos propiedades: la irreducibilidad de una forma normal y el período de una acción multiplicativa sobre ella. Este capítulo construye sus lecturas espectrales, conserva las normalizaciones de cada funcional y prepara la forma completada desarrollada en la sección~\ref{lectura:manuscrito-03}.

\subsection{Del irreducible a su reloj}
El producto y el coproducto determinan qué objetos son irreducibles. El reloj aritmético describe otra propiedad de esos objetos: el retorno de una órbita multiplicativa finita. Para un entero $n\ge2$ coprimo con diez, se define

\[
\tau_n=\min\{k\ge1:10^k\equiv1\pmod n\}.
\]
La longitud $\tau_n$ corresponde a la órbita de la unidad bajo multiplicación por diez en el grupo de unidades módulo $n$. La base decimal fija esta lectura; otra base coprima determina otro reloj. El carácter de irreducible, en cambio, pertenece al producto construido antes de escoger esta representación periódica.

\HMTresultado{Proposición}{Sincronización de relojes}{rz-num-02-1}
Si $n=\prod_p p^{a_p}$ y $\gcd(n,10)=1$, entonces

\[
\tau_n=\operatorname{mcm}_{p^{a_p}\parallel n}\tau_{p^{a_p}}.
\]
\textbf{Demostración.} Un exponente $k$ satisface $10^k\equiv1\pmod n$ si y sólo si lo satisface módulo cada potencia prima de la factorización. Esto equivale a que $k$ sea múltiplo de cada orden $\tau_{p^{a_p}}$. El menor exponente positivo común es el mínimo común múltiplo. \hfill\(\square\)

La factorización organiza, por tanto, una sincronización. Las potencias de un mismo primo deben mantenerse como tales: su período no se obtiene sustituyéndolas sin prueba por el período del primo base.

\HMTresultado{Proposición}{Cierre residual nonádico}{rz-num-02-2}
Si $\gcd(n,90)=1$, el bloque entero

\[
M_n=\frac{10^{\tau_n}-1}{n}
\]
es divisible por nueve.

\textbf{Demostración.} La definición del período hace entero a $M_n$. Como $10\equiv1\pmod9$, el numerador es divisible por nueve. El entero $n$ es invertible módulo nueve, por lo que $nM_n\equiv0\pmod9$ implica $M_n\equiv0\pmod9$. \hfill\(\square\)

Para un primo $p\notin\{2,3,5\}$, además $\tau_p\mid p-1$, pues su órbita pertenece al grupo multiplicativo de orden $p-1$. La congruencia de la proposición~\ref{rz-num-02-2} también se cumple para compuestos coprimos con noventa. Su significado es un cierre residual; la irreducibilidad continúa determinada por las aperturas del producto.

Los números $7,13,77,91,143$ tienen período decimal seis, pero sólo los dos primeros son irreducibles. La relación estructural no consiste en que un período identifique por sí solo un primo, sino en que el producto construye componentes cuya composición organiza los retornos.

\subsection{Factorización como ocupación de modos}
Sea $\mathcal P$ el conjunto obtenido evaluando los irreducibles nativos. Cada ocupación de modos es una colección $\mathbf k=(k_p)_{p\in\mathcal P}$ de enteros no negativos y soporte finito. Su evaluación multiplicativa es

\[
N(\mathbf k)=\prod_{p\in\mathcal P}p^{k_p}.
\]
La factorización única identifica estas ocupaciones con los enteros positivos. El espacio de Hilbert de las ocupaciones y el de las formas normales positivas poseen entonces una correspondencia unitaria de bases. El registro de una ocupación conserva cuántas veces interviene cada irreducible.

En el sector de ocupación total uno, \(\sum_p k_p=1\), con base \(\{e_p:p\in\mathcal P\}\), defínase

\[
H_{\mathcal P}e_p=(\log p)e_p,
\]
y, sobre las ocupaciones,

\[
H_{\mathrm{occ}}e_{\mathbf k}
=\left(\sum_pk_p\log p\right)e_{\mathbf k}.
\]
Con los dominios diagonales maximales, ambos son autoadjuntos. Si $Qe_n=ne_n$, la identificación de bases satisface

\[
U H_{\mathrm{occ}}U^*=\log Q.
\]
La igualdad es una consecuencia de la composición multiplicativa:

\[
\sum_pk_p\log p=\log N(\mathbf k).
\]
El operador no se utiliza para seleccionar retrospectivamente los primos. Recibe los irreducibles ya construidos y convierte su composición en una cantidad aditiva. Éste es el lugar preciso de la representación de ocupaciones en la genealogía del artículo.

\subsection{Una misma estructura de ocupación y dos niveles de traza}
Para $\Re s>1$, la traza sobre el sector de ocupación total uno es

\[
P(s)=\operatorname{Tr}(e^{-sH_{\mathcal P}})=\sum_{p\in\mathcal P}p^{-s}.
\]
La traza sobre todas las ocupaciones, en cambio, es

\[
Z(s)=\operatorname{Tr}(e^{-sH_{\mathrm{occ}}})
=\prod_{p\in\mathcal P}(1-p^{-s})^{-1}
=\sum_{n\ge1}n^{-s}.
\]
La última igualdad procede de la factorización única. El semiplano indicado garantiza convergencia absoluta y permite agrupar los términos. Tras la identificación de la realización entera, $Z$ coincide allí con la función zeta. No se ha utilizado una tabla de valores de zeta para definir los modos.

\HMTresultado{Proposición}{Jerarquía de momentos}{rz-num-02-3}
En $\Re s>1$, tomando la rama del logaritmo determinada por la expansión del producto convergente,

\[
\log Z(s)=\sum_{k\ge1}\frac{P(ks)}{k}.
\]
\textbf{Demostración.} La expansión absolutamente convergente $-\log(1-z)=\sum_{k\ge1}z^k/k$, aplicada a cada $z=p^{-s}$, da

\[
\log Z(s)=\sum_p\sum_{k\ge1}\frac{p^{-ks}}k.
\]
La convergencia absoluta autoriza intercambiar ambas sumas. \hfill\(\square\)

Los dos índices conservan funciones diferentes: $p$ identifica el modo irreducible y $k$ su multiplicidad de ocupación. La relación entre ambos no se reduce a una lista de constantes: produce una colección de momentos sobre el mismo conjunto de irreducibles.

\subsection{Meissel–Mertens y la eliminación de la contribución lineal}
La ampliación coordinada reúne la parte finita

\[
B_1=\gamma+\sum_p\left(\log(1-p^{-1})+p^{-1}\right).
\]
El término añadido $p^{-1}$ cancela exactamente el primer orden del logaritmo. Por tanto,

\[
B_1=\gamma-\sum_{k\ge2}\frac{P(k)}k.
\]
Aquí $\gamma$ designa la parte finita $\gamma_0$ del mismo funcional, cuya construcción y cota se desarrollan en la sección~\ref{lectura:manuscrito-04}; no es una entrada metrológica ni un decimal escogido para la evaluación. La fórmula relaciona dos operaciones sobre el funcional construido: extracción de la parte finita y cancelación del primer momento de cada modo.

\HMTresultado{Proposición}{Cota de la cola de momentos}{rz-num-02-4}
Para un entero $K\ge1$,

\[
0<\sum_{k>K}\frac{P(k)}k
\le\frac{2P(K+1)}{K+1}
\le\frac{2^{-K}}{K+1}\left(1+\frac2K\right).
\]
\textbf{Demostración.} Como $k\ge K+1$,

\[
\sum_{k>K}\frac{P(k)}k
\le\frac1{K+1}\sum_p\frac{p^{-(K+1)}}{1-p^{-1}}
\le\frac{2P(K+1)}{K+1}.
\]
Por comparación con todos los enteros mayores o iguales que dos,

\[
P(K+1)\le\sum_{n\ge2}n^{-(K+1)}
\le2^{-(K+1)}+\int_2^\infty x^{-(K+1)}\,dx.
\]
Evaluar la integral da la segunda cota. \hfill\(\square\)

La desigualdad permite convertir la suma finita en un intervalo, siempre que las evaluaciones finitas y $\gamma$ posean a su vez cotas controladas. La longitud de una expansión decimal no sustituye esta estimación.

\subsection{Exclusión residual y producto del patrón de separación dos}
Para el patrón de dos posiciones \(\{0,2\}\), sea $\nu_p$ el número de clases prohibidas módulo $p$. En $p=2$ ambas posiciones coinciden; para $p>2$, son dos clases distintas. La normalización local respecto de dos elecciones independientes produce

\[
\sigma_p=\frac{1-\nu_p/p}{(1-1/p)^2}.
\]
El factor $p=2$ vale dos, mientras que los restantes forman

\[
C_2=\prod_{p>2}\frac{1-2/p}{(1-1/p)^2}
=\prod_{p>2}\left(1-\frac1{(p-1)^2}\right).
\]
La serie singular completa del patrón es $2C_2$. Es conveniente reservar símbolos diferentes para el producto sobre $p>2$ y para el factor completo: la normalización distingue una contribución residual concreta, no una convención prescindible.

\HMTresultado{Proposición}{El mismo sistema de momentos con otro peso}{rz-num-02-5}
La expansión del producto satisface

\[
\log C_2=-\sum_{k\ge2}\frac{2^k-2}{k}\bigl(P(k)-2^{-k}\bigr).
\]
\textbf{Demostración.} Para cada $p>2$,

\[
\log(1-2/p)-2\log(1-1/p)
=-\sum_{k\ge2}\frac{2^k-2}{kp^k}.
\]
La suma empieza en $k=2$, porque el término lineal se cancela. Para \(p\ge3\),

\[
-\log\left(1-\frac1{(p-1)^2}\right)\le\frac4{3(p-1)^2}.
\]
La comparación con la serie de cuadrados inversos demuestra la finitud de la suma de valores absolutos y permite intercambiar el orden. La retirada del modo $2$ produce $P(k)-2^{-k}$. \hfill\(\square\)

La misma jerarquía $P(k)$ describe así dos operaciones diferentes: la cancelación lineal de Meissel–Mertens y la exclusión residual del patrón de dos posiciones. La diferencia está en el peso de cada momento y en la retirada del modo $2$, no en introducir dos catálogos de primos independientes.

Para un entero \(K\ge1\), sea

\[
S_K=\sum_{k=2}^K\frac{2^k-2}{k}\bigl(P(k)-2^{-k}\bigr),
\]
la ampliación coordinada da la cota

\[
0<-\log C_2-S_K
\le\frac3{K+1}\left(\frac23\right)^{K+1}\left(1+\frac3K\right).
\]
Se obtiene sustituyendo $2^k-2\le2^k$, sumando la serie geométrica para cada $p\ge3$, acotando $(1-2/p)^{-1}\le3$ y comparando la suma sobre primos con la suma sobre enteros $n\ge3$. Al exponentiar se obtiene un intervalo para $C_2$. La convergencia de este producto es un resultado distinto de la existencia de infinitos pares de primos separados por dos.

\subsection{Truncamiento por irreducibles y control del resto}
\label{ix-add:corte-primos}

La profundidad de ocupación $K$ y el límite del irreducible $Y$ definen
aproximaciones diferentes de la misma representación aritmética. Las cotas
anteriores conservan todos los irreducibles y truncan la profundidad. Ahora
se mantienen completas las contribuciones locales y se restringen los modos
producidos por el selector a $p\leq Y$, con $Y\geq2$ entero. Se definen
\[
 B_1(Y)=\gamma_0+\sum_{p\leq Y}
       \left(\log(1-p^{-1})+p^{-1}\right),\qquad
 C_2(Y)=\prod_{2<p\leq Y}\left(1-\frac1{(p-1)^2}\right).
\]
El producto vacío vale uno. La parte finita $\gamma_0$ es el mismo objeto
denotado por $\gamma$ en la fórmula de Meissel--Mertens; su evaluación con
resto se desarrolla en la sección~\ref{lectura:manuscrito-04}.

\HMTresultado{Proposición}{Encierros por corte de irreducibles}{ix-add:prime-cut-bound}
Para todo entero $Y\geq2$,
\[
 B_1(Y)-\frac1{2Y}\leq B_1\leq B_1(Y),\qquad
 \frac{Y-1}{Y}C_2(Y)\leq C_2\leq C_2(Y).
\]
\textbf{Demostración.}
Para cada $p\geq2$, la contribución que se resta de $B_1(Y)$ es
\[
 d_p=-\log(1-p^{-1})-p^{-1}
     =\sum_{k\geq2}\frac1{kp^k},\qquad
 0<d_p\leq\frac1{2p(p-1)}.
\]
La última desigualdad utiliza $1/k\leq1/2$ y la suma geométrica.
Por comparación positiva y telescopía,
\[
 0\leq B_1(Y)-B_1=\sum_{p>Y}d_p
 \leq\frac12\sum_{n>Y}\frac1{n(n-1)}=\frac1{2Y}.
\]
Cada factor de la cola de $C_2$ pertenece a $(0,1)$. Incluir todos los
enteros $n>Y$ disminuye el producto. Para $M>Y$,
\[
 \prod_{n=Y+1}^{M}\left(1-\frac1{(n-1)^2}\right)
 =\prod_{n=Y+1}^{M}\frac{n(n-2)}{(n-1)^2}
 =\frac{M(Y-1)}{Y(M-1)}.
\]
El límite es $(Y-1)/Y$. La convergencia del producto primo y la
comparación anterior dan el segundo encierro. \hfill$\square$

Al aumentar $Y$, $B_1(Y)$ y $C_2(Y)$ decrecen hacia sus valores respectivos.
Las anchuras garantizadas son $1/(2Y)$ y $C_2(Y)/Y$. Los encierros por
profundidad y por irreducible pueden intersectarse para refinar la evaluación;
cada uno conserva su parámetro y su demostración. El producto $C_2$
corresponde al sector $p>2$; el factor local del modo dos mantiene la
normalización $2C_2$ de la serie singular completa.


\subsection{Realización triádica del funcional}
El propietario de funcionales espectrales utiliza los ciclos $C_m=\mathbb Z/3^m\mathbb Z$ y su laplaciano

\[
(\Delta_mf)(j)=2f(j)-f(j+1)-f(j-1).
\]
Escribiendo \(N_m=3^m\), sean \(e_{m,n}\) los caracteres normalizados del ciclo y \(\widetilde n\) el representante de \(n\) entre \(-(N_m-1)/2\) y \((N_m-1)/2\). El operador con orientación firmada se define por

\[
D_m e_{m,n}=\frac{N_m}{\sqrt{\pi}}
\sin\left(\frac{\pi\widetilde n}{N_m}\right)e_{m,n}.
\]
Así queda fijado el signo y no solamente el cuadrado. La realización satisface

\[
D_m^2=\frac{3^{2m}}{4\pi}\Delta_m.
\]
Sea \(P_0\) el proyector ortogonal sobre las funciones constantes del ciclo. La coordenada $\pi$ pertenece a la generación anterior del núcleo HMT. Esta normalización es parte de la realización analítica y debe conservar su procedencia. La traza positiva retira el modo constante y cuenta una contribución por cada par de orientaciones opuestas:

\[
\operatorname{Tr}_m^+(F(D_m))
=\frac12\operatorname{Tr}\bigl(F(D_m)P_0^\perp\bigr).
\]
Con los modos de Fourier del ciclo, la traza de calor converge a

\[
\Theta_+(x)=\sum_{n\ge1}e^{-\pi n^2x},\qquad x>0.
\]
La prueba, desarrollada en la primera sección del capítulo siguiente, conserva una mayorante gaussiana independiente del nivel de refinamiento. La transformada de Mellin de este límite satisface exactamente

\[
\int_0^\infty\Theta_+(x)x^{s/2-1}\,dx
=\pi^{-s/2}\Gamma(s/2)Z(s),\qquad \Re s>1.
\]
Así se reúnen la construcción aditiva del espectro entero y la descomposición multiplicativa por irreducibles. La identidad entre los funcionales se demuestra; no se infiere de que dos evaluaciones numéricas sean próximas.

La inversión theta y la continuación de la función completada se demuestran en el capítulo siguiente. Las partes finitas, los coeficientes de Laurent y las clases residuales módulo tres requieren sus evaluaciones específicas; no se deducen de la sola convergencia de la traza de calor.

\subsubsection*{Defecto de transporte y componente uniforme}
La ampliación sobre pliegue de hojas proporciona una representación explícita del defecto de una vuelta. Sea \(N=9^m\), con \(m\ge1\), y sea \(T\) un operador unitario sobre una fibra de Hilbert \(\mathcal H\). En \(\mathbb C^N\) se utiliza el producto interior euclídeo usual y el vector unitario

\[
u_N=N^{-1/2}(1,\ldots,1).
\]
Se definen \(\iota_Nf=u_N\otimes f\), \(P_N=\iota_N\iota_N^*\) y \(Q_N=I-P_N\). El transporte aplica \(T\) una vez al completar el ciclo:

\[
U_N(T)=\sum_{r=1}^{N-1}|r+1\rangle\langle r|\otimes I
+|1\rangle\langle N|\otimes T.
\]
Sobre el estado uniforme actúa por

\[
U_N(T)\iota_N f
=\iota_N f+\frac{e_1}{\sqrt N}\otimes(T-I)f.
\]
El término añadido tiene una componente uniforme y otra ortogonal. La segunda, con la normalización del propietario, es

\[
\delta_N(T)=\frac{N}{\sqrt{N-1}}Q_NU_N(T)\iota_N
=\eta_N\otimes(T-I),
\]
donde

\[
\eta_N=\sqrt{\frac N{N-1}}\left(e_1-\frac{u_N}{\sqrt N}\right),
\qquad \|\eta_N\|=1,\qquad \eta_N\perp u_N.
\]
La igualdad se comprueba sustituyendo la acción anterior y aplicando \(Q_N\). Para dos transportes unitarios \(T,S\) sobre la misma fibra resulta

\[
\delta_N(T)^*\delta_N(S)=(I-T)^*(I-S).
\]
No se necesita conmutatividad: los productos aparecen en el orden indicado. La misma descomposición proporciona

\[
\delta_N(T)^*P_N\delta_N(S)=0,
\]
pues \(P_N\) anula el factor \(\eta_N\).

Estas identidades precisan la información conservada por la normalización del defecto y su posición respecto de la componente uniforme. También permiten transportar isométricamente los cierres de las imágenes entre dos longitudes de ciclo: el Gram es independiente de \(N\). Su dominio es el sector del defecto construido; no constituye por sí solo una identificación de todas las historias del TPK ni de cualquier proyección espectral posterior.

\subsection{La forma completada y su identidad de frontera}
El paso hacia el doble círculo utiliza una forma sobre un dominio de funciones de prueba. Su lado aritmético reúne las contribuciones primo–potencia, gamma y polar. La escritura con dos columnas,

\[
\mathscr I_{\mathrm{GW}}=J_+^*J_+-J_-^*J_-,
\]
expone una diferencia de formas. El productor del índice primo, los pesos orbitales, la normalización de Fourier y la completación deben figurar antes de identificarla con otra forma.

El propietario del capítulo 96 sitúa el paso central en un codificador $\Xi$ y una proyección ortogonal $P$, mediante

\[
\langle\Xi f,\Xi g\rangle=\langle J_+f,J_+g\rangle,
\qquad
\langle P\Xi f,P\Xi g\rangle=\langle J_-f,J_-g\rangle.
\]
Si estas dos identidades se obtienen de las acciones nativas declaradas, su diferencia da exactamente

\[
\mathscr I_{\mathrm{GW}}(f,g)
=\langle(I-P)\Xi f,(I-P)\Xi g\rangle.
\]
El capítulo siguiente demuestra la descomposición explícita por columnas y la consecuencia de los dos momentos. La aplicación al codificador nativo conserva como antecedente su acción sobre las funciones de prueba, la posición de sus coordenatizaciones respecto de $P$ y la prueba de ambos momentos. La ortogonalidad de $P$ explica la última igualdad, pero no determina por sí sola las dos identidades previas.

La telescopía de transporte conserva, por su parte, la forma

\[
\|x\|^2=\sum_{r=0}^{N-1}\|DA^rx\|^2+\|A^Nx\|^2,
\]
cuando $D^*D=I-A^*A$. Si sobre la hoja considerada $A=qV$, con $q=5/9$ y $V$ isométrico, el término final vale $q^{2N}\|x\|^2$ y tiende a cero. La prueba de esta extinción se conserva junto a la identidad que sitúa el estado $x$ en esa hoja.

\subsection{Círculo espectral y círculo de memoria}
El primer círculo es una coordenada de la recta espectral. Para $s=1/2+i\gamma$, con \(\gamma\in\mathbb R\), la transformación

\[
\tau(s)=\frac{s-1}{s}
=\frac{\gamma+i/2}{\gamma-i/2}
\]
pertenece al círculo unidad sin el punto \(1\) y es invertible mediante

\[
s=\frac1{1-\tau},\qquad
\gamma=\frac12\cot\frac{\theta}{2},\quad\tau=e^{i\theta}.
\]
El segundo círculo es la colección de caracteres de la memoria entera:

\[
\chi_\tau(n)=\tau^n,\qquad n\in\mathbb Z.
\]
La conexión nonádica incrementa la memoria en una unidad al completar una vuelta. El entrelazamiento relaciona una fase espectral con su acción sobre esa memoria; no identifica el cero número $j$ con la vuelta número $j$.

Si \(C\in\mathcal B(\mathcal H)\) es el operador de fase obtenido en la realización espectral, el calendario de nueve posiciones actúa sobre \(\mathbb C^9\otimes\mathcal H\) por

\[
S_{C,9}=\sum_{r=1}^8|r+1\rangle\langle r|\otimes I
+|1\rangle\langle9|\otimes C.
\]
Cada vector atraviesa una sola vez el enlace que contiene $C$ en una vuelta completa. Por tanto,

\[
S_{C,9}^9=I_9\otimes C.
\]
La identidad es exacta para cualquier $C$; cuando $C$ es unitario, $S_{C,9}$ también lo es. El factor contractivo $5/9$ pertenece a otra lectura del transporte y se distingue de esta fase unitaria.

El interés del capítulo no está en representar una matriz cualquiera sobre nueve posiciones. Está en reunir el origen del operador $C$, su dominio, el codificador anterior y la compatibilidad con los cilindros del TPK. Esos datos fijan qué se transporta y por qué esa realización es pertinente.

\subsection{Compatibilidad con la estructura conjunta del continuo}
El espacio de historias, las dos hojas, la memoria, la incidencia y el complejo no se reconstruyen como objetos independientes para cada sección. El ensamblaje correlativo previo se desarrolla en la sección~\ref{lectura:manuscrito-05}, con sus operaciones, condiciones de realización y cuadrados de naturalidad. Los operadores entero, de profundidad, de irreducibles y de alturas espectrales se mantienen tipados; una coincidencia de notación no los identifica.

La sección~\ref{lectura:manuscrito-03} reúne la identidad de frontera con su construcción terminal, las funciones de prueba, los términos de la forma completada y la convergencia. Las aproximaciones matriciales conservan sus tres matrices de forma, momento y residuo. Las cifras se comparan únicamente tras producir y congelar las salidas.

Esta organización conserva el argumento solicitado: construcción, operación, estructura resultante y reconocimiento. El manuscrito no cambia el núcleo causal para facilitar un cálculo posterior ni utiliza una lista de ceros como entrada del productor aritmético.
